\subsection{Sprint 1}

Demos início ao desenvolvimento das primeiras tarefas repassadas pelos AGES IV, em especial o login do sistema e um \textit{CRUD} básico dos questionários. O foco inicial esteve na autenticação, consolidando o fluxo de entrada no sistema e abrindo caminho para funcionalidades futuras. (\textit{aqui lembrar de citar que fomos divididos em squads})

Apesar dos esforços, apenas o login e algumas telas (login, visualização de questionários)foram concluídos de forma satisfatória nesta iteração. A complexidade do banco de dados, cuja modelagem já havia trazido dificuldades desde a Sprint 0, voltou a impactar o andamento, gerando atrasos significativos. Como consequência, apenas alguns endpoints ligados à autenticação e aos usuários ficaram prontos, enquanto o restante do escopo acabou comprometido pelo tempo limitado da sprint.

Esse cenário evidenciou a necessidade de melhorar a comunicação com a squad sobre o andamento das tarefas, evitando desalinhamentos e permitindo que os bloqueios sejam identificados e resolvidos mais cedo. Essa foi a principal lição aprendida, com impacto direto no planejamento das próximas entregas.
\\

\fbox{%
    \parbox{0.95\linewidth}{%
        \setlength{\parindent}{15pt}
            \itshape Esta sprint foi marcada principalmente por dois fatores: a sua curta duração e a alta complexidade da modelagem do banco de dados. Nossos AGES II dedicaram grande parte do tempo a essa atividade, e, somado ao pouco tempo disponível, isso acabou atrasando de forma considerável nossas entregas.

            Por outro lado, um aspecto bastante positivo foi perceber o quanto o time se manteve unido diante dessa tarefa desafiadora. Essa colaboração intensa contribuiu para fortalecer o entrosamento entre os AGES II e trouxe mais coesão ao grupo.
    }%
}
